\documentclass[../../../../uem_mei_q.tex]{subfiles}

\begin{document}
    $a$，$b$を定数とし，$-\myfraca{\sqrt{3}}{3}\leqq a\leqq\myfraca{\sqrt{3}}{3}$，%
    $-\myfraca{\sqrt{3}}{3}\leqq b\leqq\myfraca{\sqrt{3}}{3}$，$a\leqq b$ %
    であるものとする．
    \begingroup
        \setlength{\abovedisplayskip}{5pt}
        \setlength{\belowdisplayskip}{3pt}
        \begin{align*}
            x=\myfraca{ab+1}{\sqrt{a^2+1}\sqrt{b^2+1}}
        \end{align*}
    \endgroup
    により$x$を定めるとき，$x$の値が最も小さくなるのは $a=\:\text{\myboxd{　セ　}}$，%
    $b=\:\text{\myboxd{　ソ　}}$ のときであり，その時の$x$の値は \:\myboxd{　タ　}\:である．
    
    \vspace{.5\baselineskip}
    \textsf{セ，ソ，タの解答群}
    
    　\begin{tabularx}{\dimexpr\linewidth-2\zw}{@{}LLLL@{}}
        ⓪\; $0$ & ①\; $1$ & ②\; $-\myfraca{\sqrt{3}}{3}$ & ③\; $-\myfraca{1}{2}$ \\ [7pt]
        ④\; $-\myfraca{1}{3}$ & ⑤\; $-\myfraca{\sqrt{3}}{6}$ &% 
        ⑥\; $\myfraca{\sqrt{3}}{6}$ & ⑦\; $\myfraca{1}{3}$ \\ [7pt]
        ⑧\; $\myfraca{1}{2}$ & ⑨\; $\myfraca{\sqrt{3}}{3}$
    \end{tabularx}
\end{document}
